

\begin{enumerate}[label=\textbf{\arabic*)}, leftmargin=*]

    \item 
    
    É um sistema responsável por medir e quantificar determinada grandeza física e atuar sobre ela de forma a fazê-la agir da maneira desejada.

    \item
    
    São as variáveis de um sistema que se deseja medir e manter sob controle em determinado valor de referência (\textit{setpoint}). Em instrumentação, a grandeza física sob medição também é denominada \textit{mensurando}.

    \item 
    
    \begin{itemize}
        \item \textbf{Malha Aberta (\textit{Open-Loop}):} \textbf{Não há} medição ou realimentação da saída para ajustar o atuador. A ação de controle depende exclusivamente do sinal de entrada ou de um tempo pré-determinado, operando sem a confirmação de que a variável atingiu o valor desejado (ex.: ajuste manual do atuador).
        \item \textbf{Malha Fechada (\textit{Closed-Loop}):} \textbf{Há} medição constante da saída por meio de sensores (\textit{feedback}). O valor medido é comparado com a entrada de referência (\textit{setpoint}), gerando um sinal de erro que ajusta automaticamente a ação do controlador e do atuador.
    \end{itemize}

    \item 
    
    Trata-se de um sistema em malha fechada justamente porque realiza a comparação contínua entre a entrada de referência e a saída real do sistema, corrigindo desvios automaticamente.

    \item \textbf{Reguladores}
    
    Um regulador é um dispositivo projetado para controlar e manter constante uma grandeza física, como tensão elétrica, pressão de gases ou vazão de fluidos, mesmo diante de distúrbios externos.
    
    \textbf{Exemplos:} Reguladores de tensão em circuitos eletrônicos, reguladores de pressão de gás residencial (GLP) e válvulas reguladoras de pressão em processos industriais.

    \item \textbf{Modelo Matemático}
    
    Um modelo matemático é uma representação (conjunto de equações diferenciais ou funções de transferência) que descreve com fidelidade a dinâmica e a física de um sistema, permitindo prever e simular seu comportamento sob diferentes condições de operação.

    \item \textbf{Sistemas Contínuos e Discretos}
    
    \begin{itemize}
        \item \textbf{Sistemas de Dados Contínuos:} Possuem variáveis que se alteram de forma contínua ao longo do tempo $t \in \mathbb{R}$.
        \item \textbf{Sistemas de Dados Discretos:} Possuem variáveis que mudam ou são amostradas apenas em instantes específicos de tempo $k \in \mathbb{Z}$.
    \end{itemize}

\end{enumerate}